\subsubsection{Closest pair of points (sweepline)}
\begin{lstlisting}[style=mystyle, language=C++]
// TC: O(N log N)
// usa: encontrar el par de puntos mas cercanos en un plano
// idea: ordenar por x, mantener una "caja" (set ordenado por y) con puntos
//       cuya x esta a <= d del actual; evict y query acotados por d
// nota: hypot(dx,dy) = sqrt(dx^2+dy^2) sin overflow intermedio (mas seguro que sqrt(dx*dx+dy*dy))
#include <bits/stdc++.h>
using namespace std;

double closestPair(vector<Point>& pts){
	sort(pts.begin(), pts.end(), [](const Point& a, const Point& b){
		return a.x < b.x || (a.x == b.x && a.y < b.y);
	});
	double d = 1e18;
	set<pair<long long,long long>> box;
	int j = 0;
	for(int i = 0; i < (int)pts.size(); i++){
		long long xi = pts[i].x, yi = pts[i].y;
		while(j < i && (double)(xi - pts[j].x) > d){
			box.erase({pts[j].y, pts[j].x});
			j++;
		}
		long long ylow  = (long long)floor((double)yi - d);
		long long yhigh = (long long)ceil((double)yi  + d);
		for(auto it = box.lower_bound({ylow, LLONG_MIN});
			it != box.end() && it->first <= yhigh; ++it){
			double dist = hypot((double)(xi - it->second), (double)(yi - it->first));
			if(dist < d) d = dist;
		}
		box.insert({yi, xi});
	}
	return d;
}

int main(){
	vector<Point> pts = {{0,0}, {5,5}, {1,1}, {3,3}, {0,2}};
	cout << "closest distance = " << closestPair(pts) << "\n";  // ~1.41
	return 0;
}
\end{lstlisting}
